\documentclass[11pt,a4paper]{article}

\usepackage{amsmath, amssymb, amsthm}
\usepackage{physics}
\usepackage{bm}
\usepackage{siunitx}
\usepackage{booktabs}
\usepackage{graphicx}
\usepackage[margin=25mm]{geometry}
\usepackage{hyperref}

\title{Technical Note: block-VIEM.jl\\[4pt]
\large Volume Integral Equation Method with SWG/half-SWG Basis,\\
Duffy Singular Integration, and AIM-FFT Acceleration\\
for Light Scattering by Arbitrary Dielectric Particles}
\author{N.\ Moteki}
\date{Last updated: 2026-04-13 (corresponds to block-VIEM.jl v0.1.2)}

\begin{document}
\maketitle
\tableofcontents
\newpage

%======================================================================
\section{Introduction}
\label{sec:introduction}
%======================================================================

The volume integral equation method (VIEM) computes the electromagnetic
field scattered by an inhomogeneous dielectric particle of arbitrary
shape by solving an integral equation on the particle volume rather
than on its surface~\cite{SWG1984,VolakisSertel}. Compared with the
discrete dipole approximation (DDA)~\cite{Purcell1973,Draine1988,Goodman1991}
implemented in the companion package
\texttt{block-DDA\_Py}~\cite{Moteki2024}, a volumetric formulation with
conforming vector basis functions on a tetrahedral mesh offers two
practical advantages: (i)~shapes with curved or faceted surfaces are
represented without stair-casing, and (ii)~the number of degrees of
freedom per wavelength required for a given accuracy is considerably
smaller than for a cubic-lattice DDA, especially for high refractive
index contrast.

This document describes the theoretical background and the algorithms
implemented in \texttt{block-VIEM.jl} v0.1.2, covering:
\begin{enumerate}
  \item the EFVIE-D formulation with Schaubert--Wilton--Glisson (SWG)
        and half-SWG basis functions
        (Sec.~\ref{sec:formulation}),
  \item weakening the $1/R^3$ singularity to $1/R$ via integration by
        parts, including the boundary-face surface-correction terms
        $K^{B}$, $K^{C}$, $K^{D}$ required to restore charge neutrality
        (Sec.~\ref{sec:weakening}),
  \item evaluation of the remaining $1/R$ singular and near-singular
        volume integrals via the generalized Duffy transform
        (Sec.~\ref{sec:duffy}),
  \item FFT-accelerated matrix--vector products through the Adaptive
        Integral Method (AIM) with moment matching and precorrection
        (Sec.~\ref{sec:aim}),
  \item the direct and Krylov iterative solvers and the
        multi-orientation batch solve
        (Sec.~\ref{sec:solver}),
  \item the CAS-v2 compatible scattering observables computed from the
        solution (Sec.~\ref{sec:observables}),
  \item validation against Mie theory and the oblate-spheroid symmetry
        identity (Sec.~\ref{sec:validation}),
  \item the HDF5 output format for spheroid parameter sweeps, bit
        compatible with the \texttt{block-DDA\_Py} schema consumed by
        \texttt{build\_spheroid\_lut.py}
        (Sec.~\ref{sec:spheroid-sweep-io}).
\end{enumerate}

Throughout this note the electromagnetic equations are written in the
physics (BH83~\cite{Bohren1983}) time convention $\mathrm{e}^{-i\omega t}$,
so that an outgoing scalar Helmholtz Green's function carries the factor
$\mathrm{e}^{+ik_0 R}/(4\pi R)$. Absorbing materials have
$\Im(\epsilon_p) > 0$. Lengths are expressed in micrometres
(\si{\micro\metre}). Vectors are set in bold italic
($\bm{E}$, $\bm{k}$, $\bm{r}$), matrices and second-rank tensors in
upright bold ($\mathbf{A}$, $\mathbf{G}$), and unit vectors carry a hat
($\hat{\bm{e}}$, $\hat{\bm{n}}$).

%======================================================================
\section{VIEM Formulation}
\label{sec:formulation}
%======================================================================

%----------------------------------------------------------------------
\subsection{EFVIE-D strong form}
\label{sec:efvie-d}
%----------------------------------------------------------------------

Consider a non-magnetic ($\mu = \mu_0$) inhomogeneous dielectric particle
occupying the volume $\Omega$ with complex permittivity
$\epsilon_p(\bm{r})$, embedded in a homogeneous background of real
permittivity $\epsilon_{\text{bg}}$. The wavenumber in the background is
\begin{equation}
  k_0 = \omega\,\sqrt{\mu_0\,\epsilon_{\text{bg}}}
      = \frac{2\pi\,m_m}{\lambda_0},
  \label{eq:wavenumber}
\end{equation}
where $m_m = \sqrt{\epsilon_{\text{bg}}/\epsilon_0}$ is the refractive
index of the background medium.

We use the electric flux volume integral equation in the
\emph{D-form}~\cite{SWG1984}: the unknown field is the electric flux
density $\bm{D}(\bm{r}) = \epsilon_p(\bm{r})\,\bm{E}(\bm{r})$, whose
normal component is continuous across material interfaces, so the
scheme places all dielectric jumps inside the basis functions and not
in the expansion coefficients. Define the dielectric contrast
\begin{equation}
  \kappa(\bm{r})
    = \frac{\epsilon_p(\bm{r}) - \epsilon_{\text{bg}}}{\epsilon_p(\bm{r})}.
  \label{eq:kappa}
\end{equation}
The equivalent volume polarisation current is
$\bm{J}_v = -i\omega\,\kappa\,\bm{D}$, and the EFVIE-D strong form
reads~\cite{SWG1984,VolakisSertel}
\begin{equation}
  \bm{E}_{\text{inc}}(\bm{r})
    = \frac{\bm{D}(\bm{r})}{\epsilon_p(\bm{r})}
      \;-\;
      \frac{1}{\epsilon_{\text{bg}}}
      \bigl(k_0^{2} + \grad\div\bigr)
      \int_{\Omega} \kappa(\bm{r}')\,\bm{D}(\bm{r}')\,
                    G(\bm{r},\bm{r}')\;\dd^{3}r',
  \label{eq:efvie-strong}
\end{equation}
where $G(\bm{r},\bm{r}')$ is the free-space scalar Helmholtz Green's
function
\begin{equation}
  G(\bm{r},\bm{r}')
    = \frac{\mathrm{e}^{+i k_0 R}}{4\pi R},
  \qquad R = \abs{\bm{r}-\bm{r}'}.
  \label{eq:green-scalar}
\end{equation}

%----------------------------------------------------------------------
\subsection{Incident plane wave}
\label{sec:incident-wave}
%----------------------------------------------------------------------

The incident field is a monochromatic plane wave
\begin{equation}
  \bm{E}_{\text{inc}}(\bm{r})
    = \bm{E}_0\,\exp\!\bigl(+i k_0\,\hat{\bm{k}}_{\text{inc}}\cdot\bm{r}\bigr),
  \label{eq:incident-wave}
\end{equation}
with $\hat{\bm{k}}_{\text{inc}}$ the unit propagation direction and
$\bm{E}_0$ a complex polarisation vector perpendicular to
$\hat{\bm{k}}_{\text{inc}}$. For CAS-v2 observables the incident
polarisation is left-handed circular (matching
\texttt{block-DDA\_Py}~\cite{Moteki2024}):
\begin{equation}
  \bm{E}_0
    = \frac{1}{\sqrt{2}}\,\hat{\bm{e}}_{\theta,\text{inc}}
      + \frac{i}{\sqrt{2}}\,\hat{\bm{e}}_{\phi,\text{inc}},
  \label{eq:circular-pol}
\end{equation}
where $\hat{\bm{e}}_{\theta,\text{inc}},\,\hat{\bm{e}}_{\phi,\text{inc}}$
are the polar and azimuthal unit vectors of the incident direction in
the particle frame. Particle orientation is parameterised by intrinsic
ZYZ Euler angles $(\alpha,\beta,\gamma)$, as in
\texttt{block-DDA\_Py}; the laboratory-frame triad
$\{\hat{\bm{u}}_{\text{inc},L},\hat{\bm{e}}_{\theta,L},\hat{\bm{e}}_{\phi,L}\}$
is transformed into the particle frame via the inverse rotation
$\mathbf{R}(\alpha,\beta,\gamma)^{-1}$.

%----------------------------------------------------------------------
\subsection{SWG and half-SWG basis functions}
\label{sec:swg}
%----------------------------------------------------------------------

The particle volume $\Omega$ is discretised into a conforming
tetrahedral mesh $\mathcal{T}_h$ with $N_{\text{tet}}$ elements. On each
internal face $S_n$ shared by two tetrahedra $T_n^{+}$ and $T_n^{-}$, an
SWG basis function~\cite[Eqs.~(10), (12)]{SWG1984} is defined by
\begin{equation}
  \bm{f}_n(\bm{r})
  =
  \begin{cases}
    \displaystyle\frac{a_n}{3\,V_n^{+}}\,\bigl(\bm{r} - \bm{p}_n^{+}\bigr),
      & \bm{r}\in T_n^{+},\\[6pt]
    -\displaystyle\frac{a_n}{3\,V_n^{-}}\,\bigl(\bm{r} - \bm{p}_n^{-}\bigr),
      & \bm{r}\in T_n^{-},\\[6pt]
    \bm{0}, & \text{otherwise},
  \end{cases}
  \label{eq:swg-def}
\end{equation}
where $a_n$ is the face area, $V_n^{\pm}$ the volumes of the two
supporting tetrahedra, and $\bm{p}_n^{\pm}$ the free vertex (the
vertex of $T_n^{\pm}$ not on $S_n$). The normalisation in
\eqref{eq:swg-def} enforces $\bm{f}_n\!\cdot\!\hat{\bm{n}} = 1$ on
$S_n$ (oriented from $T_n^{+}$ to $T_n^{-}$) and zero on the other
five faces of $T_n^{+} \cup T_n^{-}$. The divergence is piecewise
constant,
\begin{equation}
  \div\bm{f}_n(\bm{r})
  =
  \begin{cases}
    \displaystyle +\frac{a_n}{V_n^{+}}, & \bm{r}\in T_n^{+},\\[6pt]
    \displaystyle -\frac{a_n}{V_n^{-}}, & \bm{r}\in T_n^{-},\\[6pt]
    0, & \text{otherwise}.
  \end{cases}
  \label{eq:swg-div}
\end{equation}

A naive SWG expansion that uses only internal faces implicitly imposes
$\bm{D}\!\cdot\!\hat{\bm{n}} = 0$ on the particle boundary
$\partial\Omega$, leaving the surface polarisation charge unrepresented
and producing systematic cross-section errors of order
$10$--$30\%$ on Mie tests. To restore charge neutrality, the
original~\cite{SWG1984} formulation adds one half-SWG basis function per
boundary face. On a boundary face $S_n$ with a single supporting
tetrahedron $T_n^{+}$, the half-SWG function is the first branch of
\eqref{eq:swg-def} alone:
\begin{equation}
  \bm{f}_n^{\text{half}}(\bm{r})
    = \frac{a_n}{3\,V_n^{+}}\,\bigl(\bm{r} - \bm{p}_n^{+}\bigr),
  \qquad \bm{r} \in T_n^{+},
  \label{eq:half-swg}
\end{equation}
with $\bm{f}_n\!\cdot\!\hat{\bm{n}} = 1$ on $S_n$ and zero on the other
three faces of $T_n^{+}$. For a half-SWG function the normal flux does
\emph{not} vanish on $\partial\Omega$, so it carries an induced surface
polarisation charge $\sigma = \kappa\,\bm{D}\!\cdot\!\hat{\bm{n}}$ whose
integral over $S_n$ exactly cancels the bulk bound charge of the same
basis (Sec.~\ref{sec:charge-neutrality}). The electric flux density is
expanded as
\begin{equation}
  \bm{D}(\bm{r})
    \approx \sum_{n=1}^{N_{\text{dof}}} D_n\,\bm{f}_n(\bm{r}),
  \qquad
  N_{\text{dof}} = N_{\text{int. faces}} + N_{\text{bdy. faces}}.
  \label{eq:D-expansion}
\end{equation}

%----------------------------------------------------------------------
\subsection{Galerkin weak form}
\label{sec:galerkin}
%----------------------------------------------------------------------

Testing \eqref{eq:efvie-strong} with a basis function $\bm{f}_m$ and
substituting \eqref{eq:D-expansion} yields the $N_{\text{dof}} \times
N_{\text{dof}}$ linear system
\begin{equation}
  \sum_{n=1}^{N_{\text{dof}}} Z_{mn}\,D_n = b_m,
  \qquad
  b_m = \int_{\Omega} \bm{f}_m(\bm{r})\cdot\bm{E}_{\text{inc}}(\bm{r})\,\dd^{3}r,
  \label{eq:linear-system}
\end{equation}
with the Galerkin impedance matrix element
\begin{equation}
\begin{aligned}
  Z_{mn}
    &= \int_{\Omega}
         \frac{\bm{f}_m(\bm{r})\cdot\bm{f}_n(\bm{r})}{\epsilon_p(\bm{r})}\,
         \dd^{3}r \\
    &\quad- \frac{1}{\epsilon_{\text{bg}}}
       \int_{\Omega}\bm{f}_m(\bm{r})\cdot
         \Bigl[
           \bigl(k_0^{2} + \grad\div\bigr)
           \int_{\Omega}\kappa(\bm{r}')\,\bm{f}_n(\bm{r}')\,
                        G(\bm{r},\bm{r}')\,\dd^{3}r'
         \Bigr]\dd^{3}r.
\end{aligned}
  \label{eq:Z-strong}
\end{equation}
For a homogeneous particle, $\kappa$ and $\epsilon_p$ are constants and
factor out of the inner integral, giving the compact decomposition
\begin{equation}
  Z_{mn} = \frac{1}{\epsilon_p}\,M_{mn}
           - \frac{\kappa}{\epsilon_{\text{bg}}}\,K_{mn},
  \label{eq:Z-decomposition}
\end{equation}
where $M_{mn} = \int \bm{f}_m\!\cdot\!\bm{f}_n\,\dd^{3}r$ is the mass
matrix and $K_{mn}$ is the radiation kernel defined below.

%======================================================================
\section{Weakening the Singularity via Integration by Parts}
\label{sec:weakening}
%======================================================================

As written in \eqref{eq:Z-strong}, the radiation kernel
$(k_0^{2}+\grad\div) G$ contains a hyper-singular term $1/R^{3}$ coming
from $\grad\div G$. Moving the two derivatives off $G$ onto
$\bm{f}_m$ and $\bm{f}_n$ by successive integrations by parts reduces
the singularity to the weakly singular $1/R$ kernel that Duffy
quadrature can handle. Because SWG and half-SWG functions are only
piecewise smooth and only half-SWG functions have nonzero
$\bm{f}_n\!\cdot\!\hat{\bm{n}}$ on $\partial\Omega$, the by-parts
procedure generates three boundary-face contributions in addition to
the standard bulk term.

%----------------------------------------------------------------------
\subsection{Bulk--bulk term $K^{A}$}
\label{sec:K-A}
%----------------------------------------------------------------------

Apply $\grad\div$ to the inner integral and move the inner divergence
onto $\bm{f}_n$ using the divergence theorem; do the same for the outer
divergence and $\bm{f}_m$. For internal SWG basis functions
(vanishing normal flux on $\partial\Omega$) the boundary terms are
identically zero and the resulting bulk kernel is
\begin{equation}
  K_{mn}^{A}
    = \int_{\Omega}\!\int_{\Omega}
      \bigl[
        k_0^{2}\,\bm{f}_m(\bm{r})\cdot\bm{f}_n(\bm{r}')
        - (\div\bm{f}_m)(\bm{r})\,(\div'\bm{f}_n)(\bm{r}')
      \bigr]\,
      G(\bm{r},\bm{r}')\,\dd^{3}r'\,\dd^{3}r.
  \label{eq:K-A}
\end{equation}
The integrand now contains only the weakly singular $1/R$ kernel and is
evaluated tetrahedron pair by tetrahedron pair.

%----------------------------------------------------------------------
\subsection{Surface correction terms $K^{B}$, $K^{C}$, $K^{D}$}
\label{sec:K-BCD}
%----------------------------------------------------------------------

When the source basis $\bm{f}_n$ is a half-SWG function attached to a
boundary face $S_n$, the inner integration by parts leaves a surface
term from the boundary $\partial\Omega$:
\begin{equation}
  \Phi_n(\bm{r})
    \equiv \div\!\int_{\Omega} G(\bm{r},\bm{r}')\,\bm{f}_n(\bm{r}')\,\dd^{3}r'
    = \int_{\Omega} G\,\div'\bm{f}_n\,\dd^{3}r'
      \;-\;\oint_{\partial\Omega}
          G(\bm{r},\bm{r}')\,\bigl(\bm{f}_n\!\cdot\!\hat{\bm{n}}'\bigr)\,
          \dd^{2}r'.
  \label{eq:scalar-potential}
\end{equation}
Similarly, the outer integration by parts against $\bm{f}_m$ produces a
surface term if $\bm{f}_m$ is a half-SWG function on a boundary face
$S_m$. The complete kernel is
\begin{equation}
  K_{mn} = K_{mn}^{A} + K_{mn}^{B} + K_{mn}^{C} + K_{mn}^{D},
  \label{eq:K-full}
\end{equation}
with the three boundary contributions
\begin{align}
  K_{mn}^{B}
    &= \int_{\Omega}\!\oint_{\partial\Omega}
       (\div\bm{f}_m)(\bm{r})\,\bigl(\bm{f}_n\!\cdot\!\hat{\bm{n}}'\bigr)\,
       G(\bm{r},\bm{r}')\,\dd^{2}r'\,\dd^{3}r,
  \label{eq:K-B} \\
  K_{mn}^{C}
    &= \oint_{\partial\Omega}\!\int_{\Omega}
       \bigl(\bm{f}_m\!\cdot\!\hat{\bm{n}}\bigr)\,(\div'\bm{f}_n)(\bm{r}')\,
       G(\bm{r},\bm{r}')\,\dd^{3}r'\,\dd^{2}r,
  \label{eq:K-C} \\
  K_{mn}^{D}
    &= -\oint_{\partial\Omega}\!\oint_{\partial\Omega}
       \bigl(\bm{f}_m\!\cdot\!\hat{\bm{n}}\bigr)\,
       \bigl(\bm{f}_n\!\cdot\!\hat{\bm{n}}'\bigr)\,
       G(\bm{r},\bm{r}')\,\dd^{2}r'\,\dd^{2}r.
  \label{eq:K-D}
\end{align}
By the half-SWG property $\bm{f}_n\!\cdot\!\hat{\bm{n}} = 1$ on $S_n$
and zero on the remaining boundary faces, the surface integrals in
\eqref{eq:K-B}--\eqref{eq:K-D} collapse to integrals over the specific
faces $S_m$ and $S_n$:
$K^{B}_{mn}$ is non-zero only when $n$ is a boundary face,
$K^{C}_{mn}$ only when $m$ is a boundary face, and $K^{D}_{mn}$ only
when both $m$ and $n$ are boundary faces.

%----------------------------------------------------------------------
\subsection{Charge neutrality}
\label{sec:charge-neutrality}
%----------------------------------------------------------------------

The induced bound charge associated with a single half-SWG basis on a
boundary face $S_n$ is the sum of a bulk contribution
$\rho_{\text{vol}} = -\kappa\,\div\bm{D}$ and a surface contribution
$\sigma_{\text{surf}} = \kappa\,\bm{D}\!\cdot\!\hat{\bm{n}}$. Using
\eqref{eq:swg-div} and the unit-flux property on $S_n$,
\begin{align}
  Q_{\text{vol}}
    &= -\int_{T_n^{+}} \kappa\,\div\bm{f}_n\,\dd^{3}r
     = -\kappa\,\frac{a_n}{V_n^{+}}\cdot V_n^{+}
     = -\kappa\,a_n,
  \label{eq:Q-vol} \\
  Q_{\text{surf}}
    &= \int_{S_n} \kappa\,(\bm{f}_n\!\cdot\!\hat{\bm{n}})\,\dd^{2}r
     = \kappa\,a_n.
  \label{eq:Q-surf}
\end{align}
The sum $Q_{\text{vol}} + Q_{\text{surf}} = 0$ holds exactly
\emph{only} when the surface contributions
$K^{B}, K^{C}, K^{D}$ are retained. Dropping them, as an
internal-only SWG expansion does, leaves the particle with a spurious
net bound charge $-\kappa\,a_n$ per boundary face and produces
cross-section errors of order $10$--$30\%$ on Mie tests. With the full
decomposition \eqref{eq:K-full}, block-VIEM.jl reproduces the Mie sphere
absorption cross section to $\sim 0.2\%$ at $\ell_c = 0.5$ (589~DoFs),
an order-of-magnitude improvement per DoF over cubic-lattice DDA.

%----------------------------------------------------------------------
\subsection{Mass matrix}
\label{sec:mass}
%----------------------------------------------------------------------

The first term in \eqref{eq:Z-decomposition} is the sparse mass matrix
\begin{equation}
  M_{mn}
    = \int_{T_m\cap T_n} \bm{f}_m(\bm{r})\cdot\bm{f}_n(\bm{r})\,\dd^{3}r.
  \label{eq:mass}
\end{equation}
The integrand is a degree-$2$ polynomial in each tetrahedron, so the
5-point Gauss rule on the tetrahedron (degree~3 exact) evaluates
$M_{mn}$ without quadrature error.

%======================================================================
\section{Singular and Near-Singular Volume Integration}
\label{sec:duffy}
%======================================================================

%----------------------------------------------------------------------
\subsection{Generalised vertex Duffy transform}
\label{sec:duffy-vertex}
%----------------------------------------------------------------------

The weakly singular kernel $1/R$ in \eqref{eq:K-A}--\eqref{eq:K-D}
remains integrable but standard Gauss quadrature fails when the source
and test tetrahedra share a vertex, edge, or face, or when the outer
Gauss point approaches a face of the inner element. Following
Mousavi--Sukumar~\cite{MousaviSukumar2010}, all such cases are handled
uniformly by a \emph{vertex-only} Duffy transform applied to sub-tetrahedra
obtained by projecting the singular point onto the inner element.

For $\alpha = 1$ (a $1/R$ singularity at a single vertex), the
optimised transform from the unit cube
$\bm{u}=(u_1,u_2,u_3)\in[0,1]^{3}$ to barycentric coordinates
$(\lambda_1,\lambda_2,\lambda_3,\lambda_4)$ with $\lambda_1 = 1$ at the
singular vertex is~\cite[Eq.~(1) with $\beta=1$]{MousaviSukumar2010}
\begin{equation}
  \lambda_1 = 1 - u_1,\quad
  \lambda_2 = u_1(1-u_2),\quad
  \lambda_3 = u_1 u_2 (1-u_3),\quad
  \lambda_4 = u_1 u_2 u_3.
  \label{eq:duffy-map}
\end{equation}
The face $u_1 = 0$ of the cube collapses entirely to the singular
vertex, and $R \propto u_1$ to leading order along any ray departing
from it. The total Jacobian combining \eqref{eq:duffy-map} with the
affine map from reference to physical coordinates is
\begin{equation}
  J_{D}(\bm{u}) = 6\,\abs{V_{\text{tet}}}\,u_1^{2}\,u_2.
  \label{eq:duffy-jacobian}
\end{equation}
The factor $u_1^{2}$ exactly cancels the $1/R$ singularity, leaving a
smooth integrand on the unit cube that is integrated by a tensor-product
Gauss--Legendre rule. For the $K^{A}$ kernel
\eqref{eq:K-A} the recommended order is $5\times 5\times 5$
(125~points) for self and face-adjacent pairs, and a $4\times 4\times 4$
rule for edge- and vertex-adjacent pairs; a $5$-point Gauss rule on the
source tetrahedron is sufficient for all pairs that are not near the
observation point.

%----------------------------------------------------------------------
\subsection{Pair classification}
\label{sec:pair-class}
%----------------------------------------------------------------------

Given an outer quadrature point $\bm{r}$ inside the test tetrahedron
$T_m$, the source element $T_n$ is classified and integrated according
to the minimum distance:
\begin{itemize}
  \item \emph{self} ($T_n = T_m$): the outer Gauss point is the singular
        vertex. The source tetrahedron is subdivided into four
        sub-tetrahedra with the outer Gauss point as a common vertex,
        and each is integrated by the Duffy transform
        \eqref{eq:duffy-map}.
  \item \emph{face-/edge-/vertex-adjacent}: the singular point is the
        nearest vertex, the nearest point on the shared edge, or the
        shared face. Sub-tetrahedra are generated as above around that
        singular point and integrated by the same Duffy rule.
  \item \emph{far} ($R$ larger than a chosen threshold): the inner
        integral is evaluated with a plain tetrahedron Gauss rule
        (5~points). Pairs further than the AIM near-field radius are
        never integrated directly (Sec.~\ref{sec:aim}).
\end{itemize}

%----------------------------------------------------------------------
\subsection{Surface integrals for $K^{B}, K^{C}, K^{D}$}
\label{sec:surface-int}
%----------------------------------------------------------------------

The surface contributions in Sec.~\ref{sec:K-BCD} require the
integration of $G$ over a triangle, optionally combined with an inner
tetrahedron integral ($K^{B}, K^{C}$) or an outer triangle integral
($K^{D}$). We adopt a two-dimensional version of the Duffy transform on
the reference triangle: from $(u_1, u_2)\in[0,1]^{2}$ to barycentric
coordinates $(\eta_1, \eta_2, \eta_3)$,
\begin{equation}
  \eta_1 = 1 - u_1,\quad
  \eta_2 = u_1(1-u_2),\quad
  \eta_3 = u_1 u_2,
  \label{eq:duffy-2d}
\end{equation}
with Jacobian $J_{D,2}(\bm{u}) = 2\abs{A_{\text{tri}}}\,u_1$. The
linear factor $u_1$ cancels the $1/R$ singularity at the vertex $\eta_1
= 1$, and the face $u_1 = 0$ collapses to that vertex just as in the
three-dimensional case. For the $K^{D}$ self/edge/vertex branches, the
outer triangle is first subdivided into sub-triangles at the singular
feature and the 2D Duffy rule is applied to each.

The near-singular sub-cases of $K^{B}$ and $K^{C}$ (for example, the
outer tetrahedron $T_m$ having a face coincident with the inner
boundary face $S_n$) are handled by projecting the outer Gauss point
onto the inner triangle and applying the 2D Duffy rule around that
projection.

%======================================================================
\section{FFT-Accelerated MVP: Adaptive Integral Method}
\label{sec:aim}
%======================================================================

Direct computation and storage of the dense impedance matrix $Z$ scales
as $\mathcal{O}(N_{\text{dof}}^{2})$ in memory and
$\mathcal{O}(N_{\text{dof}}^{3})$ in operations for direct factorisation,
quickly becoming prohibitive. The Adaptive Integral Method
(AIM)~\cite{ShengSong} accelerates the matrix--vector product by
representing the far-field part of the kernel on an auxiliary cubic
grid and evaluating the resulting block-Toeplitz sum with a
three-dimensional FFT, mirroring the Goodman--Draine--Flatau
construction~\cite{Goodman1991} used by
\texttt{block-DDA\_Py}.

%----------------------------------------------------------------------
\subsection{Auxiliary grid and stencil}
\label{sec:aim-grid}
%----------------------------------------------------------------------

A Cartesian grid of spacing $h_g$ is overlaid on the bounding box of
the particle. Let $\bm{g}_p$ denote the position of grid point $p$.
Each basis function $\bm{f}_n$ is associated with a stencil
$\mathcal{S}_n$ of $M_{\text{st}} = (n_{\text{st}}+1)^3$ nearby grid
points surrounding its centroid. For a moment matching order
$n_{\text{st}}$, the stencil is a cube with
$n_{\text{st}}+1$ grid points on each side.

%----------------------------------------------------------------------
\subsection{Moment matching}
\label{sec:aim-moments}
%----------------------------------------------------------------------

For each basis function $\bm{f}_n$ and each stencil point
$p\in\mathcal{S}_n$ we construct a projection weight
$W_{np}\in\mathbb{C}^{3}$ so that the multipole moments of the point
source at $p$ match those of the continuous source $\bm{f}_n$ up to
order $n_{\text{st}}$. The moments are defined with respect to the
basis centroid $\bm{c}_n$:
\begin{equation}
  \mu_n^{\bm{\alpha}}
    = \int_{\Omega}\bm{f}_n(\bm{r})\,(\bm{r}-\bm{c}_n)^{\bm{\alpha}}\,
      \dd^{3}r,
  \qquad
  \abs{\bm{\alpha}}=\alpha_x+\alpha_y+\alpha_z \le n_{\text{st}},
  \label{eq:aim-moments}
\end{equation}
and the weights $\{W_{np}\}_{p\in\mathcal{S}_n}$ are defined by the
linear constraint
\begin{equation}
  \sum_{p\in\mathcal{S}_n} W_{np}\,(\bm{g}_p - \bm{c}_n)^{\bm{\alpha}}
    = \mu_n^{\bm{\alpha}},
  \qquad \forall\,\abs{\bm{\alpha}}\le n_{\text{st}}.
  \label{eq:aim-matching}
\end{equation}
A companion set of weights matches the divergence moments
$\int (\div\bm{f}_n)(\bm{r}-\bm{c}_n)^{\bm{\alpha}}\,\dd^{3}r$ for the
scalar potential part of the kernel. The linear system
\eqref{eq:aim-matching} is small ($M_{\text{st}}$ unknowns per basis
function, $\tbinom{n_{\text{st}}+3}{3}$ equations) and is solved once,
during matrix setup, per basis function.

%----------------------------------------------------------------------
\subsection{Far-field product via FFT}
\label{sec:aim-fft}
%----------------------------------------------------------------------

With the moment matching complete, the far-field part of the
matrix--vector product $\bm{y}^{\text{far}} = K^{\text{far}}\,\bm{x}$ is
evaluated in three steps:
\begin{enumerate}
  \item \emph{Projection.} For each basis function $n$ and each stencil
        point $p$, accumulate
        $q_p = \sum_n W_{np}\,x_n$ onto the grid.
  \item \emph{Convolution.} Compute
        $Q_p = \sum_{p'} G(\bm{g}_p - \bm{g}_{p'})\,q_{p'}$ via a
        three-dimensional FFT of $q$, pointwise multiplication by the
        precomputed FFT of $G$ on the doubled grid, and an inverse FFT.
  \item \emph{Interpolation.} For each test function $m$,
        $y_m^{\text{far}} = \sum_{p\in\mathcal{S}_m} W_{mp}^{\top}\,Q_p$.
\end{enumerate}
Because $G$ depends only on the grid-point separation, its convolution
matrix is block-Toeplitz and is computed with the $3$D FFT of the same
construction used in \texttt{block-DDA\_Py}~\cite{Goodman1991}. The cost
is $\mathcal{O}(N_g \log N_g)$ per MVP, where $N_g$ is the number of
grid points. Memory for the grid Green function is also
$\mathcal{O}(N_g)$ on the doubled domain.

%----------------------------------------------------------------------
\subsection{Precorrection of near-field pairs}
\label{sec:aim-precorrection}
%----------------------------------------------------------------------

The FFT sum above is accurate only in the far field; for pairs with
$\norm{\bm{c}_m - \bm{c}_n} \le r_{\text{near}}$ it is replaced by the
exact integrals of Sec.~\ref{sec:duffy}. The construction is a
\emph{difference} of two quantities: for each near pair $(m,n)$ we
compute the exact $K_{mn}$ and subtract from it the contribution that
the FFT would have produced through the grid projection, storing the
difference in a sparse precorrection matrix $K^{\text{pc}}$. The
complete MVP is
\begin{equation}
  \bm{y} = \bigl(D_M - \tfrac{\kappa}{\epsilon_{\text{bg}}}
              (K^{\text{far}} + K^{\text{pc}})\bigr)\,\bm{x},
  \label{eq:aim-mvp}
\end{equation}
where $D_M$ is the sparse diagonal mass contribution
$M/\epsilon_p$ and $K^{\text{pc}}$ has
$\mathcal{O}(N_{\text{dof}})$ non-zero entries. The default near radius
is $r_{\text{near}} = 2 \times (\text{stencil extent})$, verified in
Phase-3 benchmarks to produce $<0.5\%$ relative error between the
AIM-accelerated and the dense matrix--vector products.

%======================================================================
\section{Solvers}
\label{sec:solver}
%======================================================================

%----------------------------------------------------------------------
\subsection{Direct and Krylov solvers}
\label{sec:direct-krylov}
%----------------------------------------------------------------------

For small problems ($N_{\text{dof}} \lesssim$ a few~thousand), the
dense assembly of $Z$ and its LU factorisation is practical and serves
as a validation reference for the AIM path. For larger problems we use
a Krylov iterative method (BiCGSTAB from
\texttt{Krylov.jl}~\cite{KrylovJl}) applied to the AIM operator
\eqref{eq:aim-mvp}. A Jacobi preconditioner formed from the diagonal of
the sparse mass and precorrection contributions suffices in our tests.

%----------------------------------------------------------------------
\subsection{Multi-orientation batch solve}
\label{sec:multi-orient}
%----------------------------------------------------------------------

Because the impedance matrix $Z$ depends only on the particle geometry
and permittivity, it can be assembled and factorised once and then
reused across many orientations: each orientation changes only the
right-hand side $\bm{b}$. The function
\texttt{solve\_cas\_v2\_orientations} assembles $Z$, computes its LU
factorisation, and then back-substitutes for every
$(\alpha_\ell,\beta_\ell,\gamma_\ell)$ in the input list, returning one
\texttt{CASv2Result} per orientation. For spheroids this is particularly
efficient: only the $\alpha = 0$ orientation needs an independent solve
for each $\beta$, and the remaining $\alpha$ values are obtained by the
analytical expansion of Sec.~\ref{sec:alpha-expansion}.

%======================================================================
\section{Scattering Observables}
\label{sec:observables}
%======================================================================

Let $D_n$ ($n = 1,\ldots,N_{\text{dof}}$) denote the solution of
\eqref{eq:linear-system} and $\bm{D}(\bm{r}) = \sum D_n\,\bm{f}_n(\bm{r})$
the reconstructed electric flux density. This section collects the
observables computed from $D_n$ that correspond to the CAS-v2 detector
protocol~\cite{Moteki2024} and can be compared directly with
\texttt{block-DDA\_Py}.

%----------------------------------------------------------------------
\subsection{Far-field scattering amplitude}
\label{sec:far-field}
%----------------------------------------------------------------------

The scattered far field in the physics convention reads
\begin{equation}
  \bm{E}^{\text{sca}}(\bm{r}) \;\sim\;
    \frac{\mathrm{e}^{+i k_0 r}}{4\pi\,r}\,\bm{F}(\hat{\bm{r}}),
  \qquad r\to\infty,
  \label{eq:far-field}
\end{equation}
where the vector scattering amplitude is obtained by transverse
projection of the Fourier transform of $\bm{D}$:
\begin{equation}
  \bm{F}(\hat{\bm{r}})
    = \frac{k_0^{2}\,\kappa}{\epsilon_{\text{bg}}}
      \bigl(\mathbf{I}_3 - \hat{\bm{r}}\otimes\hat{\bm{r}}\bigr)\cdot
      \sum_{n=1}^{N_{\text{dof}}} D_n
      \int_{\Omega}\bm{f}_n(\bm{r}')\,
                    \mathrm{e}^{-i k_0\,\hat{\bm{r}}\cdot\bm{r}'}\,
      \dd^{3}r'.
  \label{eq:F-viem}
\end{equation}
The standard scattering amplitude used in Bohren--Huffman
\cite{Bohren1983} is $\bm{f}(\hat{\bm{r}}) = \bm{F}(\hat{\bm{r}})/(4\pi)$.

%----------------------------------------------------------------------
\subsection{Absorption cross section}
\label{sec:C-abs}
%----------------------------------------------------------------------

The absorption cross section is evaluated directly from the volumetric
Joule-loss formula, which is bilinear in $\bm{D}$ and hence reduces to
a quadratic form in the expansion coefficients:
\begin{equation}
  C_{\text{abs}}
    = \frac{k_0\,\Im(\epsilon_p)}
           {\epsilon_{\text{bg}}\,\abs{\epsilon_p}^{2}\,\abs{\bm{E}_0}^{2}}\,
      \Re\!\bigl(\bm{D}^{\dagger}\mathbf{M}\,\bm{D}\bigr),
  \label{eq:C-abs}
\end{equation}
where $\mathbf{M}$ is the mass matrix \eqref{eq:mass}. For real
$\epsilon_p$ this identically vanishes.

%----------------------------------------------------------------------
\subsection{Scattering cross section by far-field integration}
\label{sec:C-sca}
%----------------------------------------------------------------------

The scattering cross section is computed as the integral of the
radiated flux density over the unit sphere,
\begin{equation}
  C_{\text{sca}}
    = \frac{1}{16\pi^{2}\,\abs{\bm{E}_0}^{2}}
      \int_{S^{2}} \abs{\bm{F}(\hat{\bm{r}})}^{2}\,\dd\Omega,
  \label{eq:C-sca}
\end{equation}
using a product Gauss--Legendre (in $\cos\theta$) times trapezoid (in
$\phi$) rule with $2 n_{\theta}^{2}$ directions. The default
$n_{\theta} = 10$ is sufficient for size parameters
$k_0\,r_{ve}\lesssim 5$; larger targets use $n_\theta\ge 20$. The
quadratic functional $\abs{\bm{F}}^{2}$ is nonnegative and insensitive
to the sign-convention subtleties that plague forward-direction
optical-theorem estimates of $C_{\text{ext}}$, so in practice we obtain
$C_{\text{ext}}$ from
\begin{equation}
  C_{\text{ext}} = C_{\text{abs}} + C_{\text{sca}}.
  \label{eq:C-ext}
\end{equation}

%----------------------------------------------------------------------
\subsection{PCAS forward amplitude}
\label{sec:S-fw}
%----------------------------------------------------------------------

The CAS-v2 forward scattering observable is defined with a circular
incident polarisation~\cite{Moteki2024}. Let
$\hat{\bm{u}}_{\text{fw}}$ be the forward scattering direction (equal
to the incident propagation direction) and
$\hat{\bm{e}}_{\theta,\text{fw}},\hat{\bm{e}}_{\phi,\text{fw}}$ the
associated polar and azimuthal unit vectors in the particle frame. The
$s$- and $p$-channel CAS observables are
\begin{align}
  S_{\text{fw},\theta}^{\text{PCAS}}\;[=S_s(0)]
    &= \frac{\sqrt{2}}{4\pi}\,
       \bm{F}(\hat{\bm{u}}_{\text{fw}})\cdot\hat{\bm{e}}_{\theta,\text{fw}},
  \label{eq:S-fw-theta} \\
  S_{\text{fw},\phi}^{\text{PCAS}}\;[=S_p(0)]
    &= -\,\frac{i\,\sqrt{2}}{4\pi}\,
       \bm{F}(\hat{\bm{u}}_{\text{fw}})\cdot\hat{\bm{e}}_{\phi,\text{fw}},
  \label{eq:S-fw-phi}
\end{align}
where the $\sqrt{2}$ and $i\sqrt{2}$ factors are inherited from the
circular polarisation decomposition \eqref{eq:circular-pol}, and the
$1/(4\pi)$ factor converts from $\bm{F}$ to the BH83-normalised
scattering amplitude $\bm{f}=\bm{F}/(4\pi)$. The PCAS forward
observable reported to the downstream CAS-v2 analyser is
\begin{equation}
  S_{\text{fw}}^{\text{PCAS}}
    = \frac{S_{\text{fw},\theta}^{\text{PCAS}}
            + S_{\text{fw},\phi}^{\text{PCAS}}}{2}.
  \label{eq:S-fw-combined}
\end{equation}

%----------------------------------------------------------------------
\subsection{OCBS backward amplitude}
\label{sec:S-bk}
%----------------------------------------------------------------------

For the exact backward scattering direction
$\hat{\bm{u}}_{\text{bk}} = -\hat{\bm{u}}_{\text{inc}}$, the laboratory
basis is
$\hat{\bm{e}}_{\theta,\text{bk}} = -\hat{\bm{e}}_{\theta,\text{inc}}$
and $\hat{\bm{e}}_{\phi,\text{bk}} = +\hat{\bm{e}}_{\phi,\text{inc}}$.
The $\theta$- and $\phi$-channel backward amplitudes are defined in
direct analogy with \eqref{eq:S-fw-theta}--\eqref{eq:S-fw-phi} at
$\hat{\bm{r}} = \hat{\bm{u}}_{\text{bk}}$, and the OCBS
(Opposite-Circular-Basis Scattering) observable~\cite{Moteki2024} is
\begin{equation}
  S_{\text{bk}}^{\text{OCBS}}
    = \frac{-\,S_{\text{bk},\theta}^{\text{OCBS}}
            + S_{\text{bk},\phi}^{\text{OCBS}}}{\sqrt{2}}.
  \label{eq:S-bk}
\end{equation}
As shown in~\cite{Moteki2024}, combining the two channels with this
sign pattern cancels the off-diagonal Mishchenko elements
$S_{12}(\pi)+S_{21}(\pi)=0$ and produces a single, orientation-stable
scalar observable.

%======================================================================
\section{Validation}
\label{sec:validation}
%======================================================================

%----------------------------------------------------------------------
\subsection{Mie sphere}
\label{sec:mie-sphere}
%----------------------------------------------------------------------

For a homogeneous sphere the exact analytical solution is the Mie
series~\cite{Bohren1983}. The reference implementation in
\texttt{test/mie\_reference.jl} provides $C_{\text{ext}}$,
$C_{\text{sca}}$, $C_{\text{abs}}$, and the full complex amplitude
matrix elements $S_{11}, S_{22}$ at $\theta=0,\pi$. The Phase-5
validation tests compare the VIEM solution to this Mie reference on a
volume-equivalent sphere with $\lambda_0 = \SI{4}{\micro\metre}$,
$m_m = 1$, $m_p = 1.5 + 0.01 i$, $r = \SI{0.5}{\micro\metre}$ and
mesh sizes $\ell_c \in\{0.30, 0.18, 0.12\}$. After the 2026-04-13
physics-convention switch (Sec.~\ref{sec:far-field}), both real and
imaginary parts of the PCAS forward amplitude
\eqref{eq:S-fw-combined} agree with the Mie reference to
$\sim 0.4\%$ at the coarsest mesh and
$\sim 0.1\%$ at $\ell_c = 0.12$. All 1704 tests in the unit-test
suite pass.

%----------------------------------------------------------------------
\subsection{Oblate spheroid $\alpha$-symmetry}
\label{sec:alpha-expansion}
%----------------------------------------------------------------------

A prolate or oblate spheroid has a continuous $C_\infty$ rotation
symmetry about its figure axis. When the orientation is parameterised
by intrinsic ZYZ Euler angles with $\gamma = 0$, rotating the first
angle $\alpha$ does not change the physical orientation of an
axisymmetric particle in the particle frame; it only rotates the
incident polarisation basis by $\alpha$ about
$\hat{\bm{u}}_{\text{inc}}$. For the CAS-v2 observables this implies
the analytical $\alpha$-expansion~\cite{Moteki2024}
\begin{align}
  S_\theta(\alpha,\beta,\gamma{=}0)
    &= A(\beta) + B(\beta)\,\mathrm{e}^{+2 i \alpha},
  \label{eq:alpha-exp-theta}\\
  S_\phi(\alpha,\beta,\gamma{=}0)
    &= A(\beta) - B(\beta)\,\mathrm{e}^{+2 i \alpha},
  \label{eq:alpha-exp-phi}
\end{align}
with
\begin{equation}
  A(\beta) = \tfrac{1}{2}\bigl[S_\theta(0,\beta,0)+S_\phi(0,\beta,0)\bigr],
  \quad
  B(\beta) = \tfrac{1}{2}\bigl[S_\theta(0,\beta,0)-S_\phi(0,\beta,0)\bigr].
  \label{eq:AB}
\end{equation}
The identity \eqref{eq:alpha-exp-theta}--\eqref{eq:alpha-exp-phi} is the
exact basis of the $\alpha$-analytical expansion that
\texttt{block-DDA\_Py} uses to reduce the orientation sweep to solves at
$\alpha = 0$ only. Because VIEM after the physics-convention switch
lives in the same time-convention as block-DDA\_Py, the sign of the
exponent is $+2i\alpha$ --- identical to block-DDA\_Py.

The benchmark script
\texttt{benchmarks/cas\_v2/spheroid\_ar3.jl} builds an oblate spheroid
with $\text{bc\_ratio}=b/c=3$, solves at $\beta = \pi/4$ for
$\alpha\in\{0, \pi/8, \pi/4, 3\pi/8, \pi/2\}$, and compares the
independently solved amplitudes to the analytical prediction
\eqref{eq:alpha-exp-theta}--\eqref{eq:alpha-exp-phi}. The maximum
relative deviation is $\sim 5\times 10^{-15}$, that is, machine
precision for double-precision arithmetic.

%----------------------------------------------------------------------
\subsection{Accuracy and computational cost vs.\ block-DDA\_Py}
\label{sec:viem-vs-dda}
%----------------------------------------------------------------------

Table~\ref{tab:viem-vs-dda-accuracy} compares the $C_{\text{abs}}$
error produced by block-VIEM.jl (half-SWG on a tetrahedral mesh) and
block-DDA\_Py (cubic-lattice DDA with the Chaumet--Rahmani dynamic
polarisability) on the same homogeneous Mie sphere,
$m_p = 1.5 + 0.01 i$, $r = \SI{1}{\micro\metre}$,
$\lambda_0 = \SI{10}{\micro\metre}$ (size parameter $x \approx 0.63$).

\begin{table}[h]
\centering
\caption{Per-DoF accuracy comparison on a Mie sphere ($m_p=1.5+0.01 i$,
$x\approx 0.63$). The VIEM solves use the half-SWG basis and Duffy
volumetric quadrature; the DDA solves use block-BiCGSTAB on the
Goodman--Draine--Flatau FFT operator.}
\label{tab:viem-vs-dda-accuracy}
\begin{tabular}{lrr}
\toprule
Method                                   & $N_{\text{dof}}$ &
$C_{\text{abs}}$ rel.\ error \\
\midrule
DDA, spacing $d = \SI{0.14}{\micro\metre}$
  & $4488$  & $1.25\%$ \\
half-SWG VIEM, $\ell_c = \SI{0.50}{\micro\metre}$
  &  $589$  & $0.21\%$ \\
half-SWG VIEM, $\ell_c = \SI{0.18}{\micro\metre}$
  & $7868$  & $0.15\%$ \\
\bottomrule
\end{tabular}
\end{table}

For this test block-VIEM.jl is roughly \textbf{10--12$\times$ more
accurate per DoF} than block-DDA\_Py: matching the $1.25\%$ DDA error
requires only $\sim 400$ half-SWG DoFs (versus $\sim 4500$ DDA dipoles),
and at the same DoF count ($\sim 8000$) the VIEM error is an order of
magnitude smaller.

There are four mutually reinforcing reasons.
\begin{enumerate}
  \item \emph{Conforming geometry.} A tetrahedral mesh represents
        curved boundaries to $O(\ell_c^{2})$, whereas a cubic lattice
        stair-cases them to $O(d)$. For Mie spheres this alone removes
        a geometric contribution to $C_{\text{abs}}$ and
        $C_{\text{sca}}$ that scales as $d$ rather than $d^{2}$.
  \item \emph{Conforming vector space.} The SWG/half-SWG basis is
        $H(\mathrm{div})$-conforming, so the normal component of
        $\bm{D}$ is automatically continuous across tetrahedron faces
        --- the physically correct boundary condition at a dielectric
        interface. Expanding $\bm{D}$ instead of
        $\bm{E}$~\cite[Sec.~II.B]{SWG1984} further ensures that the
        polarisation jump at $\partial\Omega$ is represented exactly
        by the half-SWG functions on the boundary faces, with the
        charge-neutrality identity of Sec.~\ref{sec:charge-neutrality}.
        In DDA the induced polarisation is piecewise constant at the
        dipole centres, so the interface condition is only recovered
        in the limit $d\to 0$.
  \item \emph{Linear vector basis vs point dipoles.} Within each
        tetrahedron $\bm{D}$ is represented by up to four SWG
        functions, each linear in position and with a piecewise-constant
        divergence. DDA represents the polarisation by a single point
        dipole per element, i.e.\ a delta distribution localised at
        the lattice node. The tetrahedral moment approximation is
        therefore intrinsically one order higher, and in the far field
        the error per element scales as $k_0 h$ rather than
        $(k_0 d)^{0}$.
  \item \emph{Physical polarisability.} DDA requires a dynamic
        polarisability prescription (Clausius--Mossotti plus the
        $(2/3)(ka)^{3}$ radiative correction~\cite{Draine1988}) to
        cancel leading-order interaction errors, and the choice of
        prescription feeds back into the error. VIEM solves the
        integral equation directly; $\bm{D}$ is the physical variable,
        and no such tuning is required.
\end{enumerate}

Against these accuracy advantages, block-VIEM.jl pays a higher
per-DoF setup cost. Table~\ref{tab:viem-vs-dda-cost} summarises the
asymptotic cost and memory for a matrix-free MVP:

\begin{table}[h]
\centering
\caption{Asymptotic cost model. $N$ is the number of unknowns,
$N_g$ the number of AIM or FFT grid points (typically $N_g \sim 2N$
for AIM), $n_{\text{st}}$ the AIM moment matching order, $n_{\text{it}}$
the number of Krylov iterations, and $L$ the number of particle
orientations in a batch solve.}
\label{tab:viem-vs-dda-cost}
\begin{tabular}{lll}
\toprule
Quantity                      & block-DDA\_Py
                              & block-VIEM.jl (half-SWG + AIM) \\
\midrule
Element geometry               & cubic lattice
                              & unstructured tetrahedra \\
Green kernel on grid           & analytic dyadic, $\mathcal{O}(N_g)$
                              & scalar + moment-matched dyadic,
                                $\mathcal{O}(N_g)$ \\
Near-field ``precorrection''   & none (exact on the lattice)
                              & Duffy volumetric, $\mathcal{O}(N\,n_{\text{near}})$ \\
MVP cost                       & $\mathcal{O}(N_g \log N_g)$
                              & $\mathcal{O}(N_g \log N_g + N\,n_{\text{near}})$ \\
Multi-orientation batch        & block-BiCGSTAB on $L$ RHS
                              & shared LU (small $N$) or
                                $L$ independent BiCGSTAB \\
Dominant memory                & $\mathcal{O}(N_g)$ FFT kernel
                              & $\mathcal{O}(N_g)$ FFT kernel
                                $+ \mathcal{O}(N\,n_{\text{near}})$ precorr. \\
\bottomrule
\end{tabular}
\end{table}

The central difference is that DDA has a strictly translation-invariant
interaction matrix, so a single FFT on the doubled lattice captures the
entire MVP exactly. VIEM has no such property: each basis function sees
a different discrete support of neighbouring tetrahedra, and the AIM
moment matching is only accurate for pairs whose separation is large
compared with the element size. Near pairs must therefore be computed
directly via Duffy quadrature and stored as a sparse precorrection
matrix (Sec.~\ref{sec:aim-precorrection}). For typical settings, this
near-field work dominates both the matrix assembly time and the
per-iteration MVP cost at small to moderate $N$; only for very large
meshes does the $N_g \log N_g$ FFT cost become comparable.

The overall trade-off is therefore: \textbf{VIEM spends more work per
unknown (near-field Duffy integrals and precorrection), but uses an
order of magnitude fewer unknowns to reach a given accuracy for
high-contrast particles with curved boundaries.} For sphere-like,
low-contrast problems DDA's simplicity wins on wall-clock time; for
the high-contrast iron-oxide aggregates and faceted gold nanoparticles
targeted by this project, VIEM is the clear winner and the
block-VIEM.jl $\ell_c = 0.50$ line of
Table~\ref{tab:viem-vs-dda-accuracy} is representative of production
use.

%======================================================================
\section{Spheroid Sweep HDF5 Output}
\label{sec:spheroid-sweep-io}
%======================================================================

For parameter-sweep applications the forward scattering amplitudes on a
five-dimensional grid are written to an HDF5 file whose layout is
deliberately identical to that produced by
\texttt{block-DDA\_Py}~\cite{Moteki2024}, so the downstream consumer
\texttt{PCAS\_Bayes\_APM\_Nonspherical/lut\_generation/build\_spheroid\_lut.py}
reads VIEM output without modification. The schema is documented in
\texttt{block-DDA\_Py/.claude/spheroid\_h5\_schema\_spec.md} and
summarised here for completeness.

%----------------------------------------------------------------------
\subsection{Parameter grid}
\label{sec:sweep-grid}
%----------------------------------------------------------------------

The five grid axes are stored at the root level of the file and shared
across wavelength groups:
\begin{itemize}
  \item \texttt{D\_ve\_grid} --- volume-equivalent diameter $D_{ve}$
        [\si{\micro\metre}];
  \item \texttt{RI\_real\_grid} --- real part of the particle refractive
        index, $\Re(m_p)$;
  \item \texttt{log\_AR\_grid} --- base-10 logarithm of
        $\text{bc\_ratio} = b/c$, following the
        \texttt{block-DDA\_Py} convention where $\text{bc\_ratio}>1$ is
        oblate, $<1$ is prolate, and $=1$ is a sphere;
  \item \texttt{cos\_theta\_o\_half\_grid} --- $\cos\theta_o$ on
        $[0,1]$ (half-domain, extended to $[-1,1]$ by the consumer
        using the $\theta_o\to\pi-\theta_o$ mirror symmetry);
  \item \texttt{phi\_o\_grid} --- $\phi_o$ on $[0,\pi]$ (half-domain,
        extended to $[0,2\pi]$ by the consumer using the
        $\phi_o$-periodicity $S(\theta_o,\phi_o+\pi)=S(\theta_o,\phi_o)$).
\end{itemize}
All axes must be equidistant; \texttt{write\_spheroid\_sweep\_h5}
enforces this constraint and raises an error otherwise, because the
downstream \texttt{NdimSpline\_JAX} consumer assumes uniform grids.

%----------------------------------------------------------------------
\subsection{Per-wavelength datasets}
\label{sec:sweep-data}
%----------------------------------------------------------------------

Results for each wavelength $\lambda_0$ are stored under a group
\texttt{wl\_<value>} (e.g.\ \texttt{wl\_0p638} for
$\lambda_0 = \SI{0.638}{\micro\metre}$), with the attributes
\texttt{wl\_0} and \texttt{m\_m}. Each group contains
\begin{itemize}
  \item \texttt{S\_fw\_theta\_re}, \texttt{S\_fw\_theta\_im},
        \texttt{S\_fw\_phi\_re}, \texttt{S\_fw\_phi\_im}: the real and
        imaginary parts of $S_{\text{fw},\theta}$ and
        $S_{\text{fw},\phi}$ on the shape
        $(N_{D_{ve}}, N_{RI}, N_{AR}, N_{u_{1/2}}, N_{\phi_o})$;
  \item \texttt{converged}: a boolean array of shape
        $(N_{D_{ve}}, N_{RI}, N_{AR})$ indicating successful solves.
        Non-converged cells carry \texttt{NaN+NaN$i$} in the four
        amplitude arrays across all $(N_{u_{1/2}},N_{\phi_o})$ entries.
\end{itemize}
The $\phi_o$ axis corresponds exactly to the Euler angle $\alpha$, and
the populating code applies the analytical $\alpha$-expansion
\eqref{eq:alpha-exp-theta}--\eqref{eq:alpha-exp-phi} to the $\alpha = 0$
solves, so only $N_{u_{1/2}}$ independent VIEM solves are required per
$(D_{ve},\Re(m_p),\text{AR})$ triple.

%======================================================================
\begin{thebibliography}{99}

\bibitem{SWG1984}
  D.~H.~Schaubert, D.~R.~Wilton, and A.~W.~Glisson,
  ``A tetrahedral modeling method for electromagnetic scattering by
  arbitrarily shaped inhomogeneous dielectric bodies,''
  \textit{IEEE Trans.\ Antennas Propag.} \textbf{32}(1), 77--85 (1984).

\bibitem{VolakisSertel}
  J.~L.~Volakis and K.~Sertel,
  \textit{Integral Equation Methods for Electromagnetics}
  (SciTech Publishing, 2012).

\bibitem{ShengSong}
  X.-Q.~Sheng and W.~Song,
  \textit{Essentials of Computational Electromagnetics}
  (Wiley-IEEE, 2012).

\bibitem{MousaviSukumar2010}
  S.~E.~Mousavi and N.~Sukumar,
  ``Generalized Duffy transformation for integrating vertex singularities,''
  \textit{Comput.\ Mech.} \textbf{45}(2-3), 127--140 (2010).
  \url{https://doi.org/10.1007/s00466-009-0424-1}.

\bibitem{Purcell1973}
  E.~M.~Purcell and C.~R.~Pennypacker,
  ``Scattering and absorption of light by nonspherical dielectric grains,''
  \textit{Astrophys.\ J.} \textbf{186}, 705--714 (1973).

\bibitem{Draine1988}
  B.~T.~Draine,
  ``The discrete-dipole approximation and its application to interstellar
  graphite grains,''
  \textit{Astrophys.\ J.} \textbf{333}, 848--872 (1988).

\bibitem{Goodman1991}
  J.~J.~Goodman, B.~T.~Draine, and P.~J.~Flatau,
  ``Application of fast-Fourier-transform techniques to the discrete-dipole
  approximation,''
  \textit{Opt.\ Lett.} \textbf{16}(15), 1198--1200 (1991).

\bibitem{Bohren1983}
  C.~F.~Bohren and D.~R.~Huffman,
  \textit{Absorption and Scattering of Light by Small Particles}
  (Wiley, New York, 1983).

\bibitem{Moteki2024}
  N.~Moteki and K.~Adachi,
  ``Measuring the polarized complex forward-scattering amplitudes of
  single particles in unbounded fluid flow: CAS-v2 protocol,''
  \textit{Opt.\ Express} \textbf{32}(21), 36500--36522 (2024).
  \url{https://doi.org/10.1364/OE.533776}.

\bibitem{KrylovJl}
  A.~Montoison and D.~Orban,
  ``Krylov.jl: A Julia basket of hand-picked Krylov methods,''
  \textit{J.\ Open Source Softw.} \textbf{8}(89), 5187 (2023).
  \url{https://doi.org/10.21105/joss.05187}.

\end{thebibliography}

\section*{Acknowledgment}

This document was prepared with the assistance of Claude (Anthropic).
The author assumes full responsibility for the content.

\end{document}
